\begin{center}
	\Large{\textbf{EXECUTIVE SUMMARY}}
\end{center}
\vspace{1em}
\normalsize

When patients search for medical information online, what they find is often outdated, aimed at a different clinical context, or wrong. This report describes a healthcare question-answering system that grounds each answer in a curated medical knowledge base instead of whatever a language model happened to memorise during training.

The retrieval stage combines \ac{bm25} with dense vector search and merges ranked lists with \ac{rrf}. TinyLlama~1.1B handles generation on \ac{cpu}, with no cloud inference or external \ac{api} calls. After generation, a five-signal confidence scorer (including a DeBERTa \ac{nli} check) can flag weakly supported claims before the user sees the answer.

The knowledge base draws from three sources: ChatDoctor-HealthCareMagic patient-doctor conversations, PubMedQA biomedical research abstracts, and MedMCQA postgraduate medical entrance questions. After processing with a custom RecursiveSentenceChunker, the corpus contains 505,584 documents. The chunking strategy alone accounted for a 0.292 point gain in keyword coverage, pushing the score from 0.3845 under flat character-based chunking to 0.6765 under sentence-boundary-aware chunking, measured on a 97-question evaluation set.

Three generation configurations were tested. The base TinyLlama model scored 0.3845 keyword coverage. \ac{qlora} fine-tuning over 500 steps on MedMCQA and PubMedQA training pairs raised that to 0.4634. BioMistral-7B (GGUF Q4\_K\_M\allowbreak{} quantisation) produced the best ROUGE-L score, 0.299 against TinyLlama's 0.208, but its more discursive output style resulted in lower keyword coverage. On the 36-question end-to-end benchmark, the retrieval pipeline achieved a 97.2\% hit rate.

The full system runs as a FastAPI service with a React frontend and starts with a single \texttt{make\allowbreak{} run} command. No \ac{gpu} is required.
